% ruby-tate.tex -- ruby.sty provider adapter, tate placement.
% 自己表明型 (RUBY-PASS/RUBY-FAIL をログへ出力; runner の T8 ruby ゲートが grep)。
% 検査対象（「縦ではreading frameの右にルビを配置し、
% annotationの字形はそれぞれのlocal directionへ従う」）:
%   1. ルビは読み-frame 右にあること (annotation 列 x > base 列 x)
%   2. annotation は上→下に読め、各グリフは縦字形 (直立) であること
%      = 出荷回転後の最終ページで annotation が base の行方向に沿って
%      並ぶ (y 方向に単調増加)
%   3. base 進行は通常の本文グリフと同じ行方向ピッチ (1zw)
%   4. 境界グルー: 、→ルビ の 0.5zw ペアグルーが 1 回だけ出る
%      (本文直接組との box 幅差分 0)
%   5. 連続ルビ: overlap の kern(\@tempdimb) が生きる
%   6. 状態隔離: ルビの後の和文境界が base 右端クラスで評価される
% 測定は \pdfsavepos (aux 往環) ではなく box 内寸法で行う (tate は
% 出荷回転のため最終幾何は runner の PDF 幾何ゲートが担当)。
\documentclass[pdflatex,tate,paper=a5]{jlreq}
\pagestyle{empty}
\usepackage{ruby}
\makeatletter
\newbox\ttboxa
\newbox\ttboxb
\newdimen\ttdiff
\long\def\TTEST#1#2#3#4{%
  \setbox\ttboxa=\hbox{#3}%
  \setbox\ttboxb=\hbox{#4}%
  \ttdiff=\dimexpr\wd\ttboxa-\wd\ttboxb\relax
  \ifdim\ttdiff<\z@ \ttdiff=-\ttdiff \fi
  \ifdim\ttdiff<#2\relax
    \typeout{RUBY-PASS: #1 (diff=\the\ttdiff)}%
  \else
    \typeout{RUBY-FAIL: #1 (diff=\the\ttdiff, ruby=\the\wd\ttboxa, plain=\the\wd\ttboxb)}%
  \fi}

\begin{document}
% 1. ルビ箱の「上」方向 = 内枠 (inner horizontal frame) の height。
%    tate 出荷回転 (-90°) で内枠の +y は最終ページの +x (読みframe の
%    右) に写る。よって「ルビ全体 box の height が base だけの box より
%    高い」= annotation は内枠上 = 最終右。
\setbox\ttboxa=\hbox{\ruby{漢}{かん}}%
\setbox\ttboxb=\hbox{漢}%
\ifdim\ht\ttboxa>\ht\ttboxb
  \typeout{RUBY-PASS: ann-inner-above (ht=\the\ht\ttboxa\space vs base \the\ht\ttboxb; maps to final RIGHT under -90 shipout)}%
\else
  \typeout{RUBY-FAIL: ann-inner-above (ht=\the\ht\ttboxa\space vs base \the\ht\ttboxb)}%
\fi
% 2. annotation の box が base の行方向に「並ぶ」: ルビ box の幅
%    (内枠進行方向) = base 幅 (ann が狭い場合、外枠は base 幅)。
\setbox\ttboxa=\hbox{\ruby{漢}{か}}%
\ifdim\wd\ttboxa=\dimexpr1\jlreq@jareg\relax
  \typeout{RUBY-PASS: ann-line-advance (wd=\the\wd\ttboxa)}%
\else
  \typeout{RUBY-FAIL: ann-line-advance (wd=\the\wd\ttboxa, expected \the\jlreq@jareg)}%
\fi
% 3. base の内部進行: 依存 = 2zw (tate JFM クラス幅)
\setbox\ttboxa=\hbox{\ruby{依存}{いぞん}}%
\ifdim\wd\ttboxa=\dimexpr2\jlreq@jareg\relax
  \typeout{RUBY-PASS: base-advance-2zw (wd=\the\wd\ttboxa)}%
\else
  \typeout{RUBY-FAIL: base-advance-2zw (wd=\the\wd\ttboxa, expected \the\dimexpr2\jlreq@jareg\relax)}%
\fi
% 4. 境界グルー: 、→ルビ (0.5zw) — 直接組対照との幅差 0
\TTEST{punct-left-glue-tate}{0.02pt}{漢、\ruby{字}{じ}あ}{漢、字あ}
% ルビ右→和文 (base 右端クラス)
\TTEST{right-glue-cjk-tate}{0.02pt}{漢\ruby{字}{じ}あ}{漢字あ}
% ルビ右→欧文 (xkanjiskip)
\TTEST{right-glue-latin-tate}{0.02pt}{漢\ruby{字}{じ}a}{漢字a}
% 5. 連続ルビ (overlap): 2 つ目のルビは 1 つ目の ruby@width を進む
\setbox\ttboxa=\hbox{\ruby{漢}{かんたんいそん}\ruby{字}{じ}}%
\ttdiff=\wd\ttboxa
\advance\ttdiff-\dimexpr2\jlreq@jareg+0.9\jlreq@jareg\relax
\ifdim\ttdiff<\z@ \ttdiff=-\ttdiff \fi
\ifdim\ttdiff<0.01pt
  \typeout{RUBY-PASS: consecutive-overlap-tate (wd=\the\wd\ttboxa)}%
\else
  \typeout{RUBY-FAIL: consecutive-overlap-tate (wd=\the\wd\ttboxa, expected 2.9\jlreq@jareg, residual=\the\ttdiff)}%
\fi
% 6. 状態隔離: base が句読点で終わるルビの後の和文境界 (pair 7→0)
\TTEST{state-exit-baseclass-tate}{0.02pt}{\ruby{漢、}{かん}字}{漢、字}
% 7. 本文より長いルビ (nooverlap): 外枠 = ann 幅
\rubynooverlap
\setbox\ttboxa=\hbox{\ruby{漢}{かんたんいそんですよ}}%
\ttdiff=\wd\ttboxa
\advance\ttdiff-4\jlreq@jareg
\ifdim\ttdiff<\z@ \ttdiff=-\ttdiff \fi
\ifdim\ttdiff<0.01pt
  \typeout{RUBY-PASS: long-ann-nooverlap-tate (wd=\the\wd\ttboxa)}%
\else
  \typeout{RUBY-FAIL: long-ann-nooverlap-tate (wd=\the\wd\ttboxa, expected 4\jlreq@jareg, diff=\the\ttdiff)}%
\fi
\rubyoverlap
% 8. 書体・サイズ変更
{\gtfamily\ruby{漢}{かん}ゴシック}
{\fontsize{14pt}{14pt}\selectfont\ruby{漢}{かん}大きい}
\typeout{RUBY-PASS: font-change-compiled-tate}%
% 本文実組 (PDF 幾何ゲートの対象: annotation 右・縦字形)
縦組みの本文に\ruby{依存}{いぞん}という語を埋め込む。次は、\ruby{試験}{テスト}である。
\typeout{RUBY-DONE-TATE}
\end{document}
